% Copyright (c) 2026 Eugene Joseph M. Ragasa.
% SPDX-License-Identifier: Apache-2.0
% Author: Eugene Joseph M. Ragasa
% Formal proof completed: 2026-10-19
\documentclass[11pt]{article}
\usepackage[T1]{fontenc}
\usepackage{amsmath,amssymb,booktabs,tabularx}
\usepackage[margin=1in]{geometry}
\usepackage[hidelinks]{hyperref}
\title{Proof Dossier: Relaxed Thresholds Preserve Feasibility}
\author{Eugene Joseph M. Ragasa}
\date{Formal proof completed: 2026-10-19}
\begin{document}\maketitle
\section{Identity and status}
\begin{tabularx}{\textwidth}{@{}l>{\raggedright\arraybackslash}X@{}}\toprule
Statement ID & \texttt{BG-DECISION-THM-010} \\
Status & proved; no admissions \\
Lean declaration & \path{ProjectKoios.BandGauge.Decision.thresholdFeasible_of_thresholdsRelaxed} \\
Owning module & \texttt{ProjectKoios/BandGauge/Decision/FinitePareto.lean} \\
Axioms & none \\
\bottomrule\end{tabularx}
\section{Statement}
If every relaxed budget is at least its corresponding strict budget, any decision
feasible under the strict budgets remains feasible under the relaxed budgets.
\section{Argument}
For each objective, compose the decision-to-strict-budget inequality with the
strict-to-relaxed-budget inequality.
\section{Boundary and provenance}
The theorem is monotonicity for fixed declared objective values. It provides no authority
to tune thresholds after observing results. Lean checked the exact statement; the
dossier and proof were AI-assisted.
\end{document}
